% MATH3701 Supplementary Proofs — shared preamble
% % MATH3701 Supplementary Proofs — shared preamble
% % MATH3701 Supplementary Proofs — shared preamble
% \input{proof_preamble} at the top of each proof file

\documentclass[11pt,a4paper]{article}
\usepackage{amsmath,amssymb,amsthm}
\usepackage{mathtools}
\usepackage{bm}
\usepackage{geometry}
\geometry{margin=2.5cm}
\usepackage{hyperref}
\hypersetup{colorlinks=true, linkcolor=blue, urlcolor=blue}

% ---- Math shortcuts (matching slides/preamble.tex) ----
\newcommand{\E}{\mathbb{E}}
\newcommand{\Var}{\mathrm{Var}}
\newcommand{\Cov}{\mathrm{Cov}}
\newcommand{\Prob}{\mathbb{P}}
\newcommand{\R}{\mathbb{R}}
\newcommand{\N}{\mathcal{N}}
\newcommand{\bX}{\mathbf{X}}
\newcommand{\bx}{\mathbf{x}}
\newcommand{\by}{\mathbf{y}}
\newcommand{\bY}{\mathbf{Y}}
\newcommand{\bb}{\mathbf{b}}
\newcommand{\ba}{\mathbf{a}}
\newcommand{\bbeta}{\boldsymbol{\beta}}
\newcommand{\beps}{\boldsymbol{\epsilon}}
\newcommand{\bmu}{\boldsymbol{\mu}}
\newcommand{\btheta}{\boldsymbol{\theta}}
\newcommand{\bU}{\mathbf{U}}
\newcommand{\bJ}{\mathbf{J}}
\newcommand{\bI}{\mathbf{I}}
\newcommand{\bH}{\mathbf{H}}
\newcommand{\bK}{\mathbf{K}}
\newcommand{\bL}{\mathbf{L}}
\newcommand{\bM}{\mathbf{M}}
\newcommand{\bS}{\mathbf{S}}
\newcommand{\bW}{\mathbf{W}}
\newcommand{\bz}{\mathbf{z}}
\newcommand{\boldeta}{\boldsymbol{\eta}}
\newcommand{\logit}{\mathrm{logit}}
\newcommand{\iid}{\stackrel{\mathrm{iid}}{\sim}}
\newcommand{\tr}{\mathrm{tr}}

% ---- Theorem environments ----
\theoremstyle{plain}
\newtheorem{proposition}{Proposition}
\newtheorem{lemma}{Lemma}
\newtheorem{theorem}{Theorem}
\theoremstyle{definition}
\newtheorem{remark}{Remark}

% ---- Module branding ----
\newcommand{\moduleheader}{%
  \noindent\textbf{MATH3701 Statistical Modelling}\\
  School of Mathematics, University of Leeds\\[6pt]
}
 at the top of each proof file

\documentclass[11pt,a4paper]{article}
\usepackage{amsmath,amssymb,amsthm}
\usepackage{mathtools}
\usepackage{bm}
\usepackage{geometry}
\geometry{margin=2.5cm}
\usepackage{hyperref}
\hypersetup{colorlinks=true, linkcolor=blue, urlcolor=blue}

% ---- Math shortcuts (matching slides/preamble.tex) ----
\newcommand{\E}{\mathbb{E}}
\newcommand{\Var}{\mathrm{Var}}
\newcommand{\Cov}{\mathrm{Cov}}
\newcommand{\Prob}{\mathbb{P}}
\newcommand{\R}{\mathbb{R}}
\newcommand{\N}{\mathcal{N}}
\newcommand{\bX}{\mathbf{X}}
\newcommand{\bx}{\mathbf{x}}
\newcommand{\by}{\mathbf{y}}
\newcommand{\bY}{\mathbf{Y}}
\newcommand{\bb}{\mathbf{b}}
\newcommand{\ba}{\mathbf{a}}
\newcommand{\bbeta}{\boldsymbol{\beta}}
\newcommand{\beps}{\boldsymbol{\epsilon}}
\newcommand{\bmu}{\boldsymbol{\mu}}
\newcommand{\btheta}{\boldsymbol{\theta}}
\newcommand{\bU}{\mathbf{U}}
\newcommand{\bJ}{\mathbf{J}}
\newcommand{\bI}{\mathbf{I}}
\newcommand{\bH}{\mathbf{H}}
\newcommand{\bK}{\mathbf{K}}
\newcommand{\bL}{\mathbf{L}}
\newcommand{\bM}{\mathbf{M}}
\newcommand{\bS}{\mathbf{S}}
\newcommand{\bW}{\mathbf{W}}
\newcommand{\bz}{\mathbf{z}}
\newcommand{\boldeta}{\boldsymbol{\eta}}
\newcommand{\logit}{\mathrm{logit}}
\newcommand{\iid}{\stackrel{\mathrm{iid}}{\sim}}
\newcommand{\tr}{\mathrm{tr}}

% ---- Theorem environments ----
\theoremstyle{plain}
\newtheorem{proposition}{Proposition}
\newtheorem{lemma}{Lemma}
\newtheorem{theorem}{Theorem}
\theoremstyle{definition}
\newtheorem{remark}{Remark}

% ---- Module branding ----
\newcommand{\moduleheader}{%
  \noindent\textbf{MATH3701 Statistical Modelling}\\
  School of Mathematics, University of Leeds\\[6pt]
}
 at the top of each proof file

\documentclass[11pt,a4paper]{article}
\usepackage{amsmath,amssymb,amsthm}
\usepackage{mathtools}
\usepackage{bm}
\usepackage{geometry}
\geometry{margin=2.5cm}
\usepackage{hyperref}
\hypersetup{colorlinks=true, linkcolor=blue, urlcolor=blue}

% ---- Math shortcuts (matching slides/preamble.tex) ----
\newcommand{\E}{\mathbb{E}}
\newcommand{\Var}{\mathrm{Var}}
\newcommand{\Cov}{\mathrm{Cov}}
\newcommand{\Prob}{\mathbb{P}}
\newcommand{\R}{\mathbb{R}}
\newcommand{\N}{\mathcal{N}}
\newcommand{\bX}{\mathbf{X}}
\newcommand{\bx}{\mathbf{x}}
\newcommand{\by}{\mathbf{y}}
\newcommand{\bY}{\mathbf{Y}}
\newcommand{\bb}{\mathbf{b}}
\newcommand{\ba}{\mathbf{a}}
\newcommand{\bbeta}{\boldsymbol{\beta}}
\newcommand{\beps}{\boldsymbol{\epsilon}}
\newcommand{\bmu}{\boldsymbol{\mu}}
\newcommand{\btheta}{\boldsymbol{\theta}}
\newcommand{\bU}{\mathbf{U}}
\newcommand{\bJ}{\mathbf{J}}
\newcommand{\bI}{\mathbf{I}}
\newcommand{\bH}{\mathbf{H}}
\newcommand{\bK}{\mathbf{K}}
\newcommand{\bL}{\mathbf{L}}
\newcommand{\bM}{\mathbf{M}}
\newcommand{\bS}{\mathbf{S}}
\newcommand{\bW}{\mathbf{W}}
\newcommand{\bz}{\mathbf{z}}
\newcommand{\boldeta}{\boldsymbol{\eta}}
\newcommand{\logit}{\mathrm{logit}}
\newcommand{\iid}{\stackrel{\mathrm{iid}}{\sim}}
\newcommand{\tr}{\mathrm{tr}}

% ---- Theorem environments ----
\theoremstyle{plain}
\newtheorem{proposition}{Proposition}
\newtheorem{lemma}{Lemma}
\newtheorem{theorem}{Theorem}
\theoremstyle{definition}
\newtheorem{remark}{Remark}

% ---- Module branding ----
\newcommand{\moduleheader}{%
  \noindent\textbf{MATH3701 Statistical Modelling}\\
  School of Mathematics, University of Leeds\\[6pt]
}


\begin{document}

\moduleheader

\begin{center}
  {\Large\bfseries Supplementary Proof: Moments of the Exponential Family}
\end{center}

\bigskip

\noindent This note proves that, for a distribution in the exponential
family, $\E[Y] = b'(\theta)$ and $\Var[Y] = \phi\,b''(\theta)$.
These results are used in Chapter~4 of the notes.

\section*{Setup}

A random variable $Y$ belongs to the \emph{exponential family} if its
density (or mass function) can be written
\begin{equation}
  f(y;\,\theta,\phi)
  = \exp\!\left\{\frac{y\theta - b(\theta)}{\phi} + c(y,\phi)\right\},
  \label{eq:expfam}
\end{equation}
where $\theta$ is the \emph{canonical parameter}, $\phi > 0$ is the
\emph{dispersion parameter}, $b(\theta)$ is a twice-differentiable
function, and $c(y,\phi)$ depends on $y$ and $\phi$ but not on
$\theta$.  We treat $\phi$ as known throughout.

The key tool is \emph{differentiation under the integral sign}
(Leibniz rule): for a density $f$ that integrates to~1 we may
differentiate the identity $\int f\,dy = 1$ with respect to the
parameter $\theta$.

\section*{Proof that $\E[Y] = b'(\theta)$}

Since $f(y;\theta,\phi)$ is a valid density (or mass function), we have
\[
  \int f(y;\theta,\phi)\,dy = 1.
\]
Differentiating both sides with respect to $\theta$,
\[
  \int \frac{\partial}{\partial\theta} f(y;\theta,\phi)\,dy = 0.
\]
Using the \emph{score identity}
$\partial f/\partial\theta = f \cdot \partial\log f/\partial\theta$,
this becomes
\[
  \int \frac{\partial \log f}{\partial\theta}\,f\,dy = 0,
  \qquad\text{i.e.}\qquad
  \E\!\left[\frac{\partial \log f}{\partial\theta}\right] = 0.
\]
For the exponential family \eqref{eq:expfam},
\[
  \frac{\partial \log f}{\partial\theta}
  = \frac{y - b'(\theta)}{\phi}.
\]
Hence
\[
  \E\!\left[\frac{Y - b'(\theta)}{\phi}\right] = 0
  \implies
  \E[Y] = b'(\theta). \qquad\square
\]

\section*{Proof that $\Var[Y] = \phi\,b''(\theta)$}

Differentiating $\E[\partial\log f/\partial\theta] = 0$ again with
respect to $\theta$,
\[
  \E\!\left[\frac{\partial^2 \log f}{\partial\theta^2}\right]
  + \E\!\left[\left(\frac{\partial\log f}{\partial\theta}\right)^{\!2}\right]
  = 0.
\]
This identity follows from the fact that
$\partial^2/\partial\theta^2 \int f\,dy = 0$ and that the cross-term
produces the sum of the two expectations shown.  The second term on the
left is the \emph{Fisher information}:
$\mathcal{I}(\theta) = \E[(\partial\log f/\partial\theta)^2]$.

For the exponential family,
\[
  \frac{\partial^2\log f}{\partial\theta^2} = -\frac{b''(\theta)}{\phi},
\]
so the identity gives
\[
  -\frac{b''(\theta)}{\phi} + \mathcal{I}(\theta) = 0
  \implies
  \mathcal{I}(\theta) = \frac{b''(\theta)}{\phi}.
\]
On the other hand,
\[
  \mathcal{I}(\theta)
  = \E\!\left[\left(\frac{Y-b'(\theta)}{\phi}\right)^{\!2}\right]
  = \frac{\Var[Y]}{\phi^2}.
\]
Combining,
\[
  \frac{\Var[Y]}{\phi^2} = \frac{b''(\theta)}{\phi}
  \implies
  \Var[Y] = \phi\,b''(\theta). \qquad\square
\]

\begin{remark}
  The function $V(\mu) = b''(\theta)$, expressed as a function of the
  mean $\mu = b'(\theta)$, is called the \emph{variance function}.
  For the Normal distribution $b(\theta)=\theta^2/2$ so $V(\mu)=1$;
  for the Poisson $b(\theta)=e^\theta$ so $V(\mu)=\mu$; for the
  Binomial with index $m$, $b(\theta)=m\log(1+e^\theta)$ so
  $V(\mu)=\mu(1-\mu/m)/m$.
\end{remark}

\end{document}
